\subsection{Correctness and Completeness}
\label{sec:correctness-and-completeness}
We show that the algorithm decides whether a valid assignment
exists and constructs one whenever it does.

\begin{theorem}[Complete constructive attribute synthesis]
\label{thm:complete-synthesis}
For a symbolic rewrite and metadata $m$ satisfying
Section~\ref{sec:formal-scope-and-guarantees}, including the requirement that every input
axis have extent at least two,

\begin{align}
 \operatorname{check}(m)&\iff\exists\beta:\operatorname{Valid}(m,\beta),
 \label{eq:complete-theorem}\\
 \operatorname{check}(m)&\Rightarrow
 \operatorname{Valid}(m,\operatorname{build}(m)).
 \label{eq:sound-theorem}
\end{align}
\end{theorem}
A failed \code{check} therefore means no valid assignment exists in the
stated domain. TEPL combines the feasibility check with the written
condition to determine applicability
(Section~\ref{sec:compilation-and-validation}).
Appendix~\ref{app:correctness-and-completeness} gives the proofs and the extension to
multiple terms.
